\PassOptionsToPackage{svgnames}{xcolor}
\documentclass[12pt, a4paper]{article}
%\usepackage[utf8]{inputenc}
%\usepackage{fontspec}
\usepackage{lmodern}
\usepackage[object=vectorian]{pgfornament}
\usepackage{lipsum,tikz}

\newcommand{\sectionline}{%
  \noindent
  \begin{center}
  {\color{DarkViolet}
    \resizebox{0.5\linewidth}{1ex}
    {{%
    {\begin{tikzpicture}
    \node  (C) at (0,0) {};
    \node (D) at (9,0) {};
    \path (C) to [ornament=85] (D);
    \end{tikzpicture}}}}}%
    \end{center}
  }

\newcommand{\sectionlinetwo}[2]{%
  \nointerlineskip \vspace{.5\baselineskip}\hspace{\fill}
  {\color{#1}
    \resizebox{0.5\linewidth}{2ex}
    {{%
    {\begin{tikzpicture}
    \node  (C) at (0,0) {};
    \node (D) at (9,0) {};
    \path (C) to [ornament=#2] (D);
    \end{tikzpicture}}}}}%
    \hspace{\fill}
    \par\nointerlineskip \vspace{.5\baselineskip}
  }
%\usepackage{opensans}
\usepackage{xeCJK}
% \setCJKmainfont[ItalicFont=楷体]{Source Han Serif TC}
% \setCJKsansfont{Source Han Sans TC}
% \newCJKfontfamily\sourcehansans{Source Han Sans TC}

% \usepackage{alltt}
% \usepackage{fontspec}
% \usepackage{xcolor}
% \usepackage{fancyhdr}
% \usepackage{lipsum}
% \usepackage{paracol}
% \setlength{\parindent}{0pt}
\setlength{\parskip}{5pt}

% \usepackage[a4paper,right=2cm, left=1.5cm,bottom=2cm,top=2cm, includeheadfoot, asymmetric]{geometry}
% \geometry{bindingoffset=1.5cm}

% 疊字包
% \usepackage{stackengine}

% 超鏈接包
\usepackage{hyperref}

%\usepackage{multirow}

% \pagestyle{fancy}
% \fancyhead{} % clear all fields
% \fancyfoot{} % clear all fields

%\chead{I am god}
%

%\fancyhead[LE, RO]{\thepage}
%\fancyhead[LO,CE]{From: K. Grant}
% \renewcommand\chaptermark[1]{\markboth{#1}{}}
% \renewcommand\sectionmark[1]{\markright{ \ #1}}
% \renewcommand{\headrulewidth}{0pt}
% \fancyhead[RO, RE]{\textit{\rightmark}}
% \fancyhead[LE, LO]{\textit{\leftmark}}
% \fancyfoot[LE, RO]{\thepage}

%\pagenumbering{gobble}

%\usepackage{enumitem}
%\usepackage{blindtext}
%\usepackage{ctable}

%\newcolumntype{q}[1]{>{\setlength{\parindent}{1em}}p{#1}}


%\usepackage{enumitem}
%\setlist[itemize]{leftmargin=*}

% \usepackage{endnotes}
% \makeatletter
% \def\enoteheading{\subsection*{註釋
%   \@mkboth{\notesname}{\notesname}}%
%   \mbox{}\par\vskip-2.3\baselineskip\noindent\rule{.5\textwidth}{0.4pt}\par\vskip\baselineskip}
% \makeatother

% % https://tex.stackexchange.com/questions/4959/how-do-i-indent-all-lines-of-a-paragraph-so-that-it-looks-like-a-blockquote
% \newenvironment{blockquote}{%
%   \par%
%   \medskip
%   \leftskip=2em\rightskip=0em%
%   \noindent\ignorespaces}{%
%   \par\medskip}

% \usepackage{titlesec}

% \setCJKmainfont[BoldFont=WenQuanYi Zen Hei, ItalicFont = STKaiti]{Noto Serif CJK TC}


% \titleformat{\chapter}[display]{\flushleft\Huge\bfseries\sourcehansans}{}{20pt}{\LARGE}

% \titleformat{\section}[display]{\flushleft\bfseries\sourcehansans}{}{16pt}{\large}

% \newcommand{\camilla}{\setlength{\parindent}{0pt}\bigskip\textbf{CAMILLA.}\setlength{\parindent}{2em}}
% \newcommand{\charles}{\setlength{\parindent}{0pt}\bigskip\textbf{CHARLES.}\setlength{\parindent}{2em}}
% \newcommand{\kate}{\setlength{\parindent}{0pt}\bigskip\textbf{KATE.}\setlength{\parindent}{2em}}
% \newcommand{\will}{\setlength{\parindent}{0pt}\bigskip\textbf{WILLIAM.}\setlength{\parindent}{2em}}
% \newcommand{\harry}{\setlength{\parindent}{0pt}\bigskip\textbf{HARRY.}\setlength{\parindent}{2em}}
% \newcommand{\james}{\setlength{\parindent}{0pt}\bigskip\textbf{JAMES.}\setlength{\parindent}{2em}}
% \newcommand{\evans}{\setlength{\parindent}{0pt}\bigskip\textbf{MR EVANS.}\setlength{\parindent}{2em}}
% \newcommand{\spencer}{\setlength{\parindent}{0pt}\bigskip\textbf{SPENCER.}\setlength{\parindent}{2em}}
% \newcommand{\cootsy}{\setlength{\parindent}{0pt}\bigskip\textbf{COOTSY.}\setlength{\parindent}{2em}}
% \newcommand{\jess}{\setlength{\parindent}{0pt}\bigskip\textbf{JESS.}\setlength{\parindent}{2em}}
% \newcommand{\stevens}{\setlength{\parindent}{0pt}\bigskip\textbf{MR STEVENS.}\setlength{\parindent}{2em}}
% \newcommand{\sarah}{\setlength{\parindent}{0pt}\bigskip\textbf{SARAH.}\setlength{\parindent}{2em}}
% \newcommand{\nick}{\setlength{\parindent}{0pt}\bigskip\textbf{NICK.}\setlength{\parindent}{2em}}
% \newcommand{\clive}{\setlength{\parindent}{0pt}\bigskip\textbf{CLIVE.}\setlength{\parindent}{2em}}
% \newcommand{\servant}{\setlength{\parindent}{0pt}\bigskip\textbf{SERVANT.}\setlength{\parindent}{2em}}
% \newcommand{\paul}{\setlength{\parindent}{0pt}\bigskip\textbf{PAUL.}\setlength{\parindent}{2em}}
% \newcommand{\speaker}{\setlength{\parindent}{0pt}\bigskip\textbf{SPEAKER OF THE HOUSE.}\setlength{\parindent}{2em}}
% \newcommand{\newpaper}{\setlength{\parindent}{0pt}\bigskip\textbf{FREE-NEWSPAPER WOMAN.}\setlength{\parindent}{2em}}
% \newcommand{\terry}{\setlength{\parindent}{0pt}\bigskip\textbf{TERRY.}\setlength{\parindent}{2em}}
% \newcommand{\gordon}{\setlength{\parindent}{0pt}\bigskip\textbf{SIR GORDON.}\setlength{\parindent}{2em}}
% \newcommand{\butler}{\setlength{\parindent}{0pt}\bigskip\textbf{BUTLER.}\setlength{\parindent}{2em}}
% \newcommand{\michael}{\setlength{\parindent}{0pt}\bigskip\textbf{SIR MICHAEL.}\setlength{\parindent}{2em}}
% \newcommand{\producer}{\setlength{\parindent}{0pt}\bigskip\textbf{TELEVISION PRODUCER.}\setlength{\parindent}{2em}}
% \newcommand{\bishop}{\setlength{\parindent}{0pt}\bigskip\textbf{ARCHBISHOP OF CANTERBURY.}\setlength{\parindent}{2em}}

% \newcommand{\zcamilla}{\setlength{\parindent}{0pt}\bigskip\textbf{\textsf{卡米拉：}}\setlength{\parindent}{2em}}
% \newcommand{\zcharles}{\setlength{\parindent}{0pt}\bigskip\textbf{\textsf{查爾斯：}}\setlength{\parindent}{2em}}
% \newcommand{\zkate}{\setlength{\parindent}{0pt}\bigskip\textbf{\textsf{凱特：}}\setlength{\parindent}{2em}}
% \newcommand{\zwill}{\setlength{\parindent}{0pt}\bigskip\textbf{\textsf{威廉：}}\setlength{\parindent}{2em}}
% \newcommand{\zharry}{\setlength{\parindent}{0pt}\bigskip\textbf{\textsf{哈利：}}\setlength{\parindent}{2em}}
% \newcommand{\zjames}{\setlength{\parindent}{0pt}\bigskip\textbf{\textsf{詹士：}}\setlength{\parindent}{2em}}
% \newcommand{\zevans}{\setlength{\parindent}{0pt}\bigskip\textbf{\textsf{艾文思先生：}}\setlength{\parindent}{2em}}
% \newcommand{\zspencer}{\setlength{\parindent}{0pt}\bigskip\textbf{\textsf{司奔斯：}}\setlength{\parindent}{2em}}
% \newcommand{\zcootsy}{\setlength{\parindent}{0pt}\bigskip\textbf{\textsf{庫斯：}}\setlength{\parindent}{2em}}
% \newcommand{\zjess}{\setlength{\parindent}{0pt}\bigskip\textbf{\textsf{潔斯：}}\setlength{\parindent}{2em}}
% \newcommand{\zstevens}{\setlength{\parindent}{0pt}\bigskip\textbf{\textsf{斯蒂芬思先生：}}\setlength{\parindent}{2em}}
% \newcommand{\zsarah}{\setlength{\parindent}{0pt}\bigskip\textbf{\textsf{莎拉：}}\setlength{\parindent}{2em}}
% \newcommand{\znick}{\setlength{\parindent}{0pt}\bigskip\textbf{\textsf{尼克：}}\setlength{\parindent}{2em}}
% \newcommand{\zclive}{\setlength{\parindent}{0pt}\bigskip\textbf{\textsf{克萊夫：}}\setlength{\parindent}{2em}}
% \newcommand{\zservant}{\setlength{\parindent}{0pt}\bigskip\textbf{\textsf{僕人：}}\setlength{\parindent}{2em}}
% \newcommand{\zpaul}{\setlength{\parindent}{0pt}\bigskip\textbf{\textsf{保羅：}}\setlength{\parindent}{2em}}
% \newcommand{\zspeaker}{\setlength{\parindent}{0pt}\bigskip\textbf{\textsf{議會議長：}}\setlength{\parindent}{2em}}
% \newcommand{\znewpaper}{\setlength{\parindent}{0pt}\bigskip\textbf{\textsf{派免費報紙的女人：}}\setlength{\parindent}{2em}}
% \newcommand{\zterry}{\setlength{\parindent}{0pt}\bigskip\textbf{\textsf{提瑞：}}\setlength{\parindent}{2em}}
% \newcommand{\zgordon}{\setlength{\parindent}{0pt}\bigskip\textbf{\textsf{戈登爵士：}}\setlength{\parindent}{2em}}
% \newcommand{\zbutler}{\setlength{\parindent}{0pt}\bigskip\textbf{\textsf{管家：}}\setlength{\parindent}{2em}}
% \newcommand{\zmichael}{\setlength{\parindent}{0pt}\bigskip\textbf{\textsf{麥釦爵士：}}\setlength{\parindent}{2em}}
% \newcommand{\zproducer}{\setlength{\parindent}{0pt}\bigskip\textbf{\textsf{電視製片人：}}\setlength{\parindent}{2em}}
% \newcommand{\zbishop}{\setlength{\parindent}{0pt}\bigskip\textbf{\textsf{坎特伯雷大主教：}}\setlength{\parindent}{2em}}

% \title{\textbf{King Charles III \textsf{查爾斯三世}}}
\title{\textbf{Trades' Unions and Strikes:\\Their Philosophy and Intention}}

% \author{Mike Bartlett \ \textit{著}\\高斯羽 \ \textit{譯}}
\author{T. J. Dunning}
\date{1860}

\begin{document}

     \maketitle

    %  \include{translator_note}

    %  \include{characters}

    %  \include{note_on_text}

    %  \include{prologue}

    %  \include{king_charles_iii_act_i}

    %  \include{king_charles_iii_act_ii}

    % \include{act_1}

    % \chapter{\LARGE TRADES' UNIONS AND STRIKES}
    \begin{center}
        {\LARGE TRADES' UNIONS AND STRIKES}
    \end{center}

    \vspace{2em}

    \sectionlinetwo{DarkGreen}{88}

    \vspace{2em}

    {\noindent \textsc{The}  object of this pamphlet is to obtain a hearing on the subject of
    "Trades' Unions" and "Strikes" from that class who, from their
    position in society, can only obtain a knowledge of it from the
    evidence of others. In setting forth the principle which calls Trade
    Societies into existence, it is not intended to defend all that has been
    done by Trades' Unions, any more than an advocate of religion would
    attempt to defend all that has been done in its name by religious professors. Trades' Unionists, generally, are neither writers, nor, to any
    great extent, able to employ writers in their behalf, nor, indeed, to the
    limited extent that they are able, do they at present see the necessity
    of doing so ; while, on the contrary, those opposed to them., from their
    wealth and position in society, form a great part of the very class by
    whom writers are directly and indirectly supported ; hence, while the
    press newspaper and periodical teems with denunciations against
    them, scarcely a voice, directly from themselves, is heard through the
    press in their behalf Indeed, it has been felt by impartial men to be
    the one great difficulty in the way of understanding the subject, that
    all they can get hold of on the matter except the merely local statements relative to a particular strike or dispute, which may then be
    pending is exclusively on one side. As far as our limited means and
    power will enable us, it is intended by this pamphlet to supply that
    deficiency, by a statement direct from a body of working men, by their
    accredited agent, of their views on the subject.}

    The principal argument against Trades' Unions is professed to be
    derived from the alleged ignorance of the working classes of Political
    Economy. As this term is used by many without having a precise
    idea of what it is, and by different persons in a different sense, we will
    commence by very briefly stating the meaning which we attach to
    it: Political Economy is the science which treats of the principle of exchange in its operation in the production of wealth, and its operation
    and influence on civilised society.

    Man is said to be an animal in contradistinction from all others
    that makes exchanges. And it is by the principle of exchange that all
    the transactions of civilised life are carried on, by which, indeed,
    civilised society exists and is held together. But if man be an animal
    that makes exchanges, he undoubtedly retains, in common with the
    carnivora to which he in part belongs, instincts of a predatory
    character, and might be defined to be an animal that, besides making
    exchanges, also steals. The act of exchange embraces every con

    ceivable form in which the subject can be treated, and, indeed, is that
out of which the science itself is evolved. It is only necessary, therefore, by analysing this act, to illustrate what we understand by the
term Political Economy. By the act of exchange, the four things
essential to a civilised community are supposed, without which, indeed,
it could not exist, namely. Labour Capital the right of Individual
Appropriation, or Property and Security. Labour, to produce the
thing exchanged ; Capital, in the accumulated labour, which results
in the article exchanged ; the recognition of the right of Individual
Appropriation, or property, by which the act of exchange takes place,
as distinguished from the act of taking or keeping possession of the
article by force ; and Security, which ensures the peaceful continuation
of a succession of exchanges to all who have anything to dispose of :
the result of the whole, in civilised life, being the four great
classes into which society is divided. Labour yielding wages
giving support to all who labour (using the term in its most
extensive sense) ; capital giving profit the revenue of those who
employ capital. The right of individual appropriation, or property, as
applied to the land, which, without remark, we take as we find
it,
gives
rent, which gives revenue to landlords. And security supposes taxes,
which support those who defend the country, administer its laws, and
-s, tax-receivers in general. This last strongly illustrates what
is
very
"^
little noticed,
if
at all, by political economists namely, the predatory
instincts of mankind, the omission of Avhich often leads to very great
'
error in the application of rules of political economy, and to none
more than to their application to Trades' Unions. And yet these
predatory instincts are most unmistakably indicated and provided
against by this element of exchange, security, in the absence of Avhich
it
certainly could not go on, as in communities
it
requires laws with
strong force to support those who administer them and enforce
their execution, without which the community could not exist ; and
in separate nations, armies and fleets for defence, and, as a
stronger illustration, for attack. On these four elements of the
principle of exchange, without either of which
it
could not exist,
it
will be seen the whole system of civilised society rests. We have thus
particularised each element, and its respective result, to show that
those who live by labour are, by right, entitled to their revenue by the
simple action of the principle of exchange equally with the other three
great divisions of society, as opposed to the confused idea of wages
being derived from the relation of master and servant, involving a sort
of benevolent patronage on the one side, and obligation on the other.
And
it
will be our object to show, that
it
is
only by trade combinations
that the free exercise of this principle can be secured to working men

    In section \nameref{marker} we defined...
    
    \section{test} \label{marker}
    
    % In section \nameref{marker} we defined...

\end{document}
